% GENERATED FILE. Edit canonical lessons, then run npm run generate:books.
% canonical-source-hash: fnv1a64:b4e7ad309bb46a1a
% canonical-lessons: SA-C435-kapardakah, SA-C435-niskah, SA-C435-dinarah, SA-C435-vinimayah, SA-C435-bhandagaram

\chapter{Old Coins and Storehouses}
\label{ch:sa-old-coins-and-storehouses}
\hlchaptermodality{435}

\begin{chapteropening}
\textbf{By the end of this chapter:} \emph{I can talk about the money of former times.}

Complete the last lesson of chapter 435: I can talk about the money of former times.
\end{chapteropening}

\section[\texorpdfstring{\sk{कपर्दकः} (kapardakaḥ)}{कपर्दकः (kapardakaḥ)}]{\sk{कपर्दकः} (kapardakaḥ) — a cowrie shell (as money)}

\label{lesson:SA-C435-kapardakah}



\begin{quote}
Before the new one: say the Sanskrit for a deposit, then the Sanskrit for pays back.
\end{quote}

\subsection*{You'll want to know: \sk{कपर्दकः}}
\textbf{\sk{कपर्दकः}} — \emph{kapardakaḥ} — \textquotedblleft{}a cowrie shell (as money)\textquotedblright{}.

\begin{quote}
\sk{पूर्वं} \sk{जनाः} \sk{कपर्दकैः} \sk{धान्यं} \sk{क्रीणन्ति} \sk{स्म।}
\end{quote}

\begin{quote}
\emph{pūrvaṃ janāḥ kapardakaiḥ dhānyaṃ krīṇanti sma.}
\end{quote}

\begin{quote}
Long ago people used to buy grain with cowrie shells.
\end{quote}

A masculine noun.

\begin{cousinweb}
\textbf{\sk{कपर्दकः}} (\emph{kapardakaḥ}) is built on \emph{kaparda}, \textquotedblleft{}a cowrie shell\textquotedblright{}, with the ending \emph{-ka}; MW gives \emph{kapardaka} as the cowrie shell, used as a small coin, but no derivation for \emph{kaparda}, so its origin is uncertain (MW).
\end{cousinweb}

\begin{grammarlens}[title={What you've built}]
One more word for getting around.
\end{grammarlens}

\subsection*{Your turn}
\begin{itemize}
\raggedright
  \item \emph{Say it:} \emph{kapardakaḥ}
  \item \emph{Say it:} \emph{kapardakaḥ}, once more
  \item \emph{Say it:} \emph{pūrvaṃ janāḥ kapardakaiḥ dhānyaṃ krīṇanti sma.} — Long ago people used to buy grain with cowrie shells.
  \item \emph{Recall:} say the Sanskrit for a deposit, then the Sanskrit for pays back, then say \emph{kapardakaḥ} again
\end{itemize}

\begin{tcolorbox}[breakable,colback=teal!4,colframe=teal!35!black,title={Before you move on}]
What does \sk{कपर्दकः} mean? (\textquotedblleft{}A cowrie shell (as money)\textquotedblright{}.) Say it once more.
\end{tcolorbox}
\section[\texorpdfstring{\sk{निष्कः} (niṣkaḥ)}{निष्कः (niṣkaḥ)}]{\sk{निष्कः} (niṣkaḥ) — a gold piece}

\label{lesson:SA-C435-niskah}



\begin{quote}
Before the new one: say the Sanskrit for pays back, then the Sanskrit for a cowrie shell.
\end{quote}

\subsection*{You'll want to know: \sk{निष्कः}}
\textbf{\sk{निष्कः}} — \emph{niṣkaḥ} — \textquotedblleft{}a gold piece\textquotedblright{}.

\begin{quote}
\sk{राजा} \sk{कवये} \sk{शतं} \sk{निष्कान्} \sk{ददाति।}
\end{quote}

\begin{quote}
\emph{rājā kavaye śataṃ niṣkān dadāti.}
\end{quote}

\begin{quote}
The king gives the poet a hundred gold pieces.
\end{quote}

A masculine noun.

\begin{cousinweb}
\textbf{\sk{निष्कः}} (\emph{niṣkaḥ}) is the word MW gives for a gold ornament for the neck and a coin or weight of gold; MW gives no derivation, so its origin is uncertain (MW).
\end{cousinweb}

\begin{grammarlens}[title={What you've built}]
One more word for getting around.
\end{grammarlens}

\subsection*{Your turn}
\begin{itemize}
\raggedright
  \item \emph{Say it:} \emph{niṣkaḥ}
  \item \emph{Say it:} \emph{niṣkaḥ}, once more
  \item \emph{Say it:} \emph{rājā kavaye śataṃ niṣkān dadāti.} — The king gives the poet a hundred gold pieces.
  \item \emph{Recall:} say the Sanskrit for pays back, then the Sanskrit for a cowrie shell, then say \emph{niṣkaḥ} again
\end{itemize}

\begin{tcolorbox}[breakable,colback=teal!4,colframe=teal!35!black,title={Before you move on}]
What does \sk{निष्कः} mean? (\textquotedblleft{}A gold piece\textquotedblright{}.) Say it once more.
\end{tcolorbox}
\section[\texorpdfstring{\sk{दीनारः} (dīnāraḥ)}{दीनारः (dīnāraḥ)}]{\sk{दीनारः} (dīnāraḥ) — a dinar (a gold coin)}

\label{lesson:SA-C435-dinarah}



\begin{quote}
Before the new one: say the Sanskrit for a cowrie shell, then the Sanskrit for a gold piece.
\end{quote}

\subsection*{You'll want to know: \sk{दीनारः}}
\textbf{\sk{दीनारः}} — \emph{dīnāraḥ} — \textquotedblleft{}a dinar (a gold coin)\textquotedblright{}.

\begin{quote}
\sk{धनिकस्य} \sk{स्यूते} \sk{दीनाराः} \sk{सन्ति।}
\end{quote}

\begin{quote}
\emph{dhanikasya syūte dīnārāḥ santi.}
\end{quote}

\begin{quote}
There are dinars in the rich man's bag.
\end{quote}

A masculine noun.

\begin{cousinweb}
\textbf{\sk{दीनारः}} (\emph{dīnāraḥ}) is borrowed from Greek \emph{dēnarion}, the Roman \emph{denarius}, as MW derives it; MW gives \emph{dīnāra} as a gold ornament and a gold coin (MW).
\end{cousinweb}

\begin{grammarlens}[title={What you've built}]
One more word for getting around.
\end{grammarlens}

\subsection*{Your turn}
\begin{itemize}
\raggedright
  \item \emph{Say it:} \emph{dīnāraḥ}
  \item \emph{Say it:} \emph{dīnāraḥ}, once more
  \item \emph{Say it:} \emph{dhanikasya syūte dīnārāḥ santi.} — There are dinars in the rich man's bag.
  \item \emph{Recall:} say the Sanskrit for a cowrie shell, then the Sanskrit for a gold piece, then say \emph{dīnāraḥ} again
\end{itemize}

\begin{tcolorbox}[breakable,colback=teal!4,colframe=teal!35!black,title={Before you move on}]
What does \sk{दीनारः} mean? (\textquotedblleft{}A dinar (a gold coin)\textquotedblright{}.) Say it once more.
\end{tcolorbox}
\section[\texorpdfstring{\sk{विनिमयः} (vinimayaḥ)}{विनिमयः (vinimayaḥ)}]{\sk{विनिमयः} (vinimayaḥ) — barter; exchange}

\label{lesson:SA-C435-vinimayah}



\begin{quote}
Before the new one: say the Sanskrit for a gold piece, then the Sanskrit for a dinar.
\end{quote}

\subsection*{You'll want to know: \sk{विनिमयः}}
\textbf{\sk{विनिमयः}} — \emph{vinimayaḥ} — \textquotedblleft{}barter; exchange\textquotedblright{}.

\begin{quote}
\sk{पूर्वं} \sk{ग्रामेषु} \sk{धान्यस्य} \sk{विनिमयः} \sk{भवति} \sk{स्म।}
\end{quote}

\begin{quote}
\emph{pūrvaṃ grāmeṣu dhānyasya vinimayaḥ bhavati sma.}
\end{quote}

\begin{quote}
Long ago grain used to be bartered in the villages.
\end{quote}

A masculine noun.

\begin{cousinweb}
\textbf{\sk{विनिमयः}} (\emph{vinimayaḥ}) is built from \emph{vi}, \emph{ni} and the root √me, \textquotedblleft{}to exchange\textquotedblright{}; MW gives \emph{vinimaya} as exchange, barter (MW).
\end{cousinweb}

\begin{grammarlens}[title={What you've built}]
One more word for getting around.
\end{grammarlens}

\subsection*{Your turn}
\begin{itemize}
\raggedright
  \item \emph{Say it:} \emph{vinimayaḥ}
  \item \emph{Say it:} \emph{vinimayaḥ}, once more
  \item \emph{Say it:} \emph{pūrvaṃ grāmeṣu dhānyasya vinimayaḥ bhavati sma.} — Long ago grain used to be bartered in the villages.
  \item \emph{Recall:} say the Sanskrit for a gold piece, then the Sanskrit for a dinar, then say \emph{vinimayaḥ} again
\end{itemize}

\begin{tcolorbox}[breakable,colback=teal!4,colframe=teal!35!black,title={Before you move on}]
What does \sk{विनिमयः} mean? (\textquotedblleft{}Barter; exchange\textquotedblright{}.) Say it once more.
\end{tcolorbox}
\section[\texorpdfstring{\sk{भाण्डागारम्} (bhāṇḍāgāram)}{भाण्डागारम् (bhāṇḍāgāram)}]{\sk{भाण्डागारम्} (bhāṇḍāgāram) — a storehouse}

\label{lesson:SA-C435-bhandagaram}



\begin{quote}
Before the new one: say the Sanskrit for a dinar, then the Sanskrit for barter.
\end{quote}

\subsection*{You'll want to know: \sk{भाण्डागारम्}}
\textbf{\sk{भाण्डागारम्}} — \emph{bhāṇḍāgāram} — \textquotedblleft{}a storehouse\textquotedblright{}.

\begin{quote}
\sk{नगरस्य} \sk{भाण्डागारे} \sk{बहु} \sk{धान्यम्} \sk{अस्ति।}
\end{quote}

\begin{quote}
\emph{nagarasya bhāṇḍāgāre bahu dhānyam asti.}
\end{quote}

\begin{quote}
There is a lot of grain in the city's storehouse.
\end{quote}

A neuter noun.

\begin{cousinweb}
\textbf{\sk{भाण्डागारम्}} (\emph{bhāṇḍāgāram}) is built from \emph{bhāṇḍa}, \textquotedblleft{}a vessel, wares\textquotedblright{}, and \emph{āgāra}, \textquotedblleft{}a house, a room\textquotedblright{}; MW gives the compound as a treasury, a storehouse (MW).
\end{cousinweb}

\begin{grammarlens}[title={What you've built}]
One more word for getting around.
\end{grammarlens}

\subsection*{Your turn}
\begin{itemize}
\raggedright
  \item \emph{Say it:} \emph{bhāṇḍāgāram}
  \item \emph{Say it:} \emph{bhāṇḍāgāram}, once more
  \item \emph{Say it:} \emph{nagarasya bhāṇḍāgāre bahu dhānyam asti.} — There is a lot of grain in the city's storehouse.
  \item \emph{Recall:} say the Sanskrit for a dinar, then the Sanskrit for barter, then say \emph{bhāṇḍāgāram} again
\end{itemize}

\begin{tcolorbox}[breakable,colback=teal!4,colframe=teal!35!black,title={Before you move on}]
What does \sk{भाण्डागारम्} mean? (\textquotedblleft{}A storehouse\textquotedblright{}.) Say it once more.
\end{tcolorbox}
